%\clearpage

\section {Introduction}
\label{sec:intro}

In this technical note, we report the results of a 3-flavour $\bar{\nu}_{\mu}$-disappearance analysis,
performed by the {\em VALOR} group,        
on the combined Run 5+6 dataset, using only data taken in antineutrino beam mode (\TotalPot{} POT).
This dataset was collected in two run periods:
May - June 2014 (Run 5, \RunVPot{} POT in antineutrino beam mode) and November 2014 - June 2015 
(Run 6, \RunVIPot{} POT in antineutrino beam mode).
This analysis builds on our $\nu_{\mu}$-disappearance analyses presented in Refs.\ \cite{OurRun1DisappAnalysis}, 
\cite{OurRun12DisappAnalysis}, \cite{OurRun12DisappAnalysis3Flavour}, \cite{OurRun123DisappAnalysis3Flavour},
\cite{OurRun1234aDisappAnalysis3Flavour} and \cite{OurRun1234DisappAnalysis3Flavour}, and adapts these analyses  
to antineutrino beam mode data.

One of the objectives of this analysis is to determine whether the oscillation parameters for $\bar{\nu}_{\mu}$ disappearance are different 
from those for $\nu_{\mu}$ disappearance. For this reason, we use two separate sets of oscillation parameters in this analysis. The first set is
$\sin^{2}\overline{\theta}_{23}$, $\sin^{2}\overline{\theta}_{12}$, $\sin^{2}\overline{\theta}_{13}$, $|\overline{\Delta m}^{2}_{32}|$ 
(normal hierarchy) or $|\overline{\Delta m}^{2}_{31}|$ (inverted hierarchy)
, $\overline{\Delta m}^{2}_{21}$ and $\delta_{CP}$, and these are 
applied to the $\bar{\nu}_{\mu}$, $\bar{\nu}_{e}$ and oscillated $\bar{\nu}_{e}$ interaction modes. The second set of oscillation parameters 
is $\sin^{2}\theta_{23}$, $\sin^{2}\theta_{12}$, $\sin^{2}\theta_{13}$, $|\Delta m^{2}_{32}|$ 
(normal hierarchy) or $|\Delta m^{2}_{31}|$ (inverted hierarchy), 
$\Delta m^{2}_{21}$ and $\delta_{CP}$, which are applied to the $\nu_{\mu}$, $\nu_{e}$ and oscillated $\nu_{e}$
interaction modes. All these parameters are fixed in the fits of the Super K data except $\sin^{2}\overline{\theta}_{23}$ and $|\overline{\Delta m}^{2}_{32}|$ 
(normal hierarchy) or $|\overline{\Delta m}^{2}_{31}|$ (inverted hierarchy).

These oscillation parameters are determined by fitting the predicted reconstructed energy spectrum of 1 $\mu$-like ring 
events at SuperK to the observed spectrum. The method of calculation of this predicted spectrum is presented in Sec.~\ref{sec:spectrum_prediction}.
This analysis uses flux tuning based on recent NA61 data \cite{FluxTuning}\cite{NA61T2K}. 
It also uses an indirect extrapolation from ND280 to SuperK, which is implemented via a fit of flux and cross-section parameters to
ND280 measurements of the muon momentum-angle distributions of $\bar{\nu}_\mu$ and  $\nu_\mu$
CC0$\pi$-like, CC1$\pi$-like and CC other events whose 
vertices are in FGD1\cite{BANFF2015}.
Updated nominal predictions of the reconstructed energy spectrum of single $\mu$-like ring events in SuperK are presented for Run 5+6
in Sec.~\ref{sec:spectrum_prediction}

The oscillation analysis is performed using neutrino oscillation probabilities computed in a 3-flavour
framework including constant-density matter effects. 
Details of our oscillation probability calculations, estimates of the numerical accuracy of our
calculations and comparisons with calculations used by the alternative T2K oscillation analyses 
are given in Appendices A and B of Ref. \cite{OurRun123DisappAnalysis3Flavour}.

As in our previous analyses, the oscillation fit is made using the likelihood-ratio method \cite{StatAnalysisCowan}.
The mathematical details of this method are given in Sec.\ 3.2 in Ref.\ \cite{OurRun1DisappAnalysis}.
We consider a total of \NSyst{} energy scale, neutrino flux, neutrino cross-section,
efficiency, initial-state nuclear environment and final-state rescattering systematics (nuisance parameters).
In this Run 5+6 analysis, we allow all these nuisance parameters to float.
The nuisance parameters and their effects 
on the extrapolated single $\mu$-like ring event reconstructed energy spectrum are described in Sec.~\ref{sec:systematics}.
The oscillation fitting method is presented in detail in Sec.~\ref{sec:oscfit_method},
while the chosen goodness-of-fit test, suitable for the current statistics, is described in Sec.~\ref{sec:gof_test}.
The performance of the fitter was tested in detail for the Run 5+6 fit using ensembles of toy MC experiments, and
a summary of the results of these tests is presented in Sec.~\ref{sec:validation}. 

The 2D confidence regions in $\sin^{2}\overline{\theta}_{23}$ and $|\overline{\Delta m}^{2}_{32}|$ 
or $|\overline{\Delta m}^{2}_{31}|$ 
are constructed using the constant-$\Delta \chi^{2}$
%~ constructed using the Feldman-Cousins 
method. 
Systematic errors are included in the oscillation contours using the profiling method.
%Systematic errors are included in the oscillation contours using the 
%Cousins-Highland method (Bayesian integration) 
%~ \cite{CousinsHighland,Kendall}.
%%Confidence contours constructed with the constant-$\Delta\chi^{2}$ method are also shown for comparison.
%Both of these methods are summarized in Sec.~\ref{sec:contour_construction}, and
%the exact methodology used for incorporating systematic errors is explained in detail.

The results of the measurements of $\sin^{2}\overline{\theta}_{23}$ and $|\overline{\Delta m}^{2}_{32}|$ 
or $|\overline{\Delta m}^{2}_{31}|$ 
using the Run 5+6 dataset are given in Sec.~\ref{sec:results_bestfit} of this note, 
while the goodness-of-fit tests are presented in Sec.~\ref{sec:results_gof}.
The best-fit spectra are shown in Sec.~\ref{sec:results_bestfit} for each of the normal and inverted mass hierarchy assumptions for the 
Run 5+6 dataset.
%~ The Feldman-Cousins and constant-$\Delta\chi^{2}$ confidence regions are shown in Sec.~\ref{sec:results_contours}. 
The confidence regions are shown in Sec.~\ref{sec:results_dchisq_contours}. 
%~ Comparisons of our results obtained between different methods for constructing confidence regions are presented in Sec.~\ref{sec:results_contour_comparisons}.
In that section, we also show comparisons of our T2K Run 5+6 results with results from MINOS
as well as with the expected T2K sensitivity for an integrated exposure of \TotalPot{} POT.

Finally, a summary is given in Sec.~\ref{sec:summary}.  

\clearpage
